% 中文论文 LaTeX 模板 (CTeX)
% 适用于中文核心期刊、学位论文
\documentclass[UTF8, a4paper, 12pt]{ctexart}

% 页面设置
\usepackage{geometry}
\geometry{
    top=2.54cm,
    bottom=2.54cm,
    left=3.17cm,
    right=3.17cm,
    headheight=10mm,
    footskip=10mm
}

% 数学与图表
\usepackage{amsmath, amssymb, amsfonts}
\usepackage{graphicx}
\usepackage{booktabs}
\usepackage{longtable}
\usepackage{hyperref}
\usepackage{cleveref}
\usepackage{listings}
\usepackage{xcolor}

% 代码高亮配置
\lstset{
    basicstyle=\ttfamily\small,
    numbers=left,
    numberstyle=\tiny\color{gray},
    frame=shadowbox,
    escapeinside={\%}{\%},
    extendedchars=false,
    language=Python,
}

% 标题设置
\ctexset{
    section = {
        format = \raggedright\Large\bfseries,
        number = \Chinese{section}
    },
    subsection = {
        format = \raggedright\large\bfseries,
    }
}

\begin{document}

% 标题
\title{基于深度学习的金融文档理解方法研究}
\author{
    姓名\qquad 学号\qquad 导师\qquad 专业\\
    某某大学\quad 某某学院
}

\date{\today}

\maketitle

\begin{abstract}
本文针对金融文档理解任务，提出了一种基于预训练语言模型的金融文档分析方法。实验表明，该方法在金融文档分类、情感分析等任务上取得了显著的性能提升。
\end{abstract}

\keywords{金融文档理解; 预训练语言模型; 深度学习; 自然语言处理}

\section{引言}

\section{相关工作}

\section{方法}

\section{实验}

\section{结论}

\section*{参考文献}

\begin{thebibliography}{99}
    \bibitem{ref1} Author. Title. Conference, Year.
\end{thebibliography}

\end{document}
ibliography{references}
